\section{Discussion and limitations}

\subsection{Discussion: coupling to The Polymathic Singularity}

The environmental extension predicts changes in payoff, not a replacement of specialists by a superior biological type. Its useful empirical unit is a task-specific interaction: performance = person factors + environment factors + person \ensuremath{\times} environment + resource and history terms. A civilization-level claim requires repeated evidence across settings rather than extrapolation from one successful human–AI collaboration.

Doshi and Hauser's experiment found that generative-AI ideas could improve evaluated individual story creativity while reducing collective diversity \cite{ref21}. This supplies a concrete example of gains and losses occurring at different scales. It does not demonstrate a special advantage for an exoregulatory phenotype, historical selection reversal, or a completed phase transition.

The \ensuremath{\kappa} conversion ratio in the earlier papers depends on an unobserved denominator, “latent integrative capacity.” A stronger operational approach measures verified output, quality, cost, time, and completion separately; any latent-capacity estimate requires an explicit measurement model. New tools can support regulation, increase distraction, or both. The practical empirical question is which coupling conditions improve sustained functioning and durable work, for whom, and at what cost.


\subsection{Limitations}

The full chain has not been demonstrated. This narrative synthesis does not establish completeness of the literature, pooled effects, publication-bias correction, or causal identification. Some pivotal studies are small or preparation-specific, and source anchors are not interchangeable with replications. Human intervention evidence does not establish the proposed mediator. Brain molecular state is difficult to measure longitudinally in living infants, creating a substantial translational gap.

RCB, PRI as an integrated operator, the broad RCCX constitutional cluster, and the Dąbrowski-to-biology mapping remain unvalidated constructs or syntheses. Their flexibility creates a risk of post hoc accommodation; preregistered exclusions, alternative models, stopping rules, and independent replication are necessary. The framework does not explain every developmental phenotype and does not infer individual etiology from group averages. It does not support reducing clinical needs to environmental mismatch or romanticizing physiological impairment.

Historical claims about industrial institutions and future labor demand are narrower here than in the conceptual source. They remain motivating interpretations until operationalized. Bibliographic verification establishes the identity and relevance of a cited study, not a forensic audit of its raw data. Journal-specific formatting, external domain review, and an author-approved disclosure statement remain submission steps.
